\documentclass[a4paper,12pt]{ctexart}
\usepackage{xcolor}
\usepackage{graphicx}
\usepackage[top=1.8cm, bottom=1.8cm, left=1.8cm, right=1.8cm]{geometry}
\usepackage{array}
\usepackage{multirow}
\usepackage{hyperref}
\hypersetup{colorlinks=true,linkcolor=blue!70!black}
\linespread{1.35}

\newcommand{\ruby}[2]{#1\textsuperscript{#2}}
\newcommand{\shihai}[2]{%
  \vspace{0.35cm}
  \noindent
  \fcolorbox{orange!60!black}{orange!8}{%
    \begin{minipage}{0.915\textwidth}
    \textbf{\textcolor{orange!70!black}{史海拾遗：}#1} \\[0.1cm]
    {\small #2}
    \end{minipage}%
  }
  \vspace{0.35cm}
}

\title{\textcolor{blue!70!black}{\Huge\textbf{欧洲古代史}}}
\author{}
\date{}

\begin{document}
\maketitle
\thispagestyle{empty}

\tableofcontents
\newpage

\section{\textcolor{orange!80!black}{第一课\quad 查理曼与欧洲的"分裂基因"}}

\subsection*{一个问题：欧洲为什么没能统一？}

我们在中国古代史里看到过一个反复出现的规律：统一→分裂→再统一→再分裂→再统一。秦统一了六国，汉亡后是三国两晋南北朝的大分裂，隋唐重新统一，五代十国再次分裂，北宋再统一……你可以把它叫做"\textbf{分久必合}"——在整个中国古代史的漫长跨度里，统一是常态，分裂是插曲。

而欧洲走了一条完全相反的路。自公元843年之后，欧洲再也没有被统一过——至今已经将近一千二百年了。拿破仑短暂地接近过，希特勒疯狂地尝试过，但最终都失败了。\textbf{"合久必分"}——这是中国的故事，在欧洲则更像是"\textbf{分了就再也没有合起来}"。

这不是某个皇帝英明或愚蠢的结果。当你比较中国和欧洲的地图时，你会发现\u2014\u2014\textbf{地理本身在悄悄决定历史走向}。

\subsection*{中国为什么能"合"？}

中国的地理条件天然有利于统一：中原是一块巨大的平原——黄河和长江之间没有不可逾越的山脉，从西安到洛阳到开封到南京，军队可以长驱直入。这种平坦的地形意味着\u2014\u2014\textbf{统一的成本低，分裂的维护成本高}。一个小政权在大平原上很难长期独立存在——你的邻居随时可以平推过来。

刚统一的秦朝做了另一件关键的事：用统一的文字（书同文）和道路系统（驰道/直道）把这片巨大的平原编织成了一个紧密的整体。哪怕后来在政治上分裂（比如三国、南北朝），\textbf{所有人的脑子里认的还是同一个文字、同一套经典、同一个"天下"}。所以当武力重新统一的时候，文化上不需要从头整合。

\subsection*{欧洲为什么"合"不了？}

现在看欧洲地图。欧洲的海岸线极其破碎，内部被阿尔卑斯山、比利牛斯山、喀尔巴阡山脉切成了互不连通的小块。一条宽阔的莱茵河就能让两岸发展出完全不同的方言和文化。在这种地形下，建立一个统一的帝国需要跨越高山、大河、海峡——成本极高；而分裂后的小政权依靠天险可以轻松防守——\textbf{分裂的维护成本低，统一的管理成本高}。

这不仅仅是地形的问题。公元800年的圣诞节，法兰克国王\textbf{查理曼}（查理大帝）在罗马的圣彼得大教堂被教皇加冕为"\textbf{罗马人的皇帝}"。他建立了一个涵盖今天法国、德国、意大利北部、低地国家的庞大帝国——这是西罗马帝国灭亡300多年后，西欧第一次出现如此规模的统一政权。

但查理曼死后仅仅29年，公元843年，他的三个孙子在\textbf{凡尔登}签了一份协议——把祖父的帝国分成了三块：

\begin{center}
\begin{tabular}{|c|c|}
\hline
\textbf{分给谁} & \textbf{大致范围} \\
\hline
长子洛泰尔 & 中间狭长地带（今荷兰、阿尔萨斯、洛林、意大利北部） \\
\hline
次子路易 & 东法兰克（今德国的雏形） \\
\hline
三子查理 & 西法兰克（今法国的雏形） \\
\hline
\end{tabular}
\end{center}

\shihai{分家产的逻辑——一种制度为什么坚持了一千多年？}{
  法兰克人——和后来的绝大多数中世纪欧洲王国——实行的继承规则和古代中国完全不同。

  中国在周代就确立了\textbf{嫡长子继承制}：父亲死后，正妻所生的长子继承全部家业和头衔。这套制度确保了权力和土地不同分割——代代传下去，王国还是完整的。汉武帝的"推恩令"后来故意打破了这条规则来削弱诸侯，但那是另一个故事。

  法兰克人的规矩是\textbf{诸子平分}。这套制度来自日耳曼部落的传统：父亲的财产是所有儿子共有的遗产，每个儿子都应该分到一份。王国和地产一样——是可分的财产。查理曼的帝国之所以在843年被他的孙子们一分为三，不是因为他们特别不团结，而是因为按照他们的继承习惯法，帝国本来就该这么分。

  但这个看似"公平"的制度，产生了一个在数学上很可怕的长期后果：分一次、再分一次、再分一次——王国的领地越分越碎。法国的卡佩王朝一开始直接控制的领地小得可怜——"法兰西岛"——就是巴黎周边小小一块，比好几个大贵族的领地还小。中国家族分家产也是诸子平分，但政治权力不分——皇位只有一个继承人。欧洲的规则把政治权力也一起分了。不同的继承制度，在一千年的时间尺度下，积累出了截然不同的政治版图。
}

三兄弟签完协议之后，各自的后代又继续在自己的地盘上按照同样的规则往下分。此后的欧洲再也没有回到过查理曼时代的统一程度。分裂不是一次性的偶然事件——是一个\textbf{内置在继承规则里的持续过程}。

\subsection*{本课要点}

\begin{center}
\begin{tabular}{|c|c|p{7cm}|}
\hline
\textbf{事件} & \textbf{年代} & \textbf{意义} \\
\hline
查理曼加冕 & 800年 & 西罗马灭亡后西欧最大的统一政权 \\
\hline
凡尔登条约 & 843年 & 帝国一分为三→法、德、意雏形；欧洲"再也合不起来"的起点 \\
\hline
\end{tabular}
\end{center}

\vspace{0.3cm}
{\large\textcolor{orange!80!black}{\textbf{想一想}}}
\begin{enumerate}
  \item 中国的"分久必合"和欧洲的"分了就合不起来"——地理、文字和继承制度各起了什么作用？如果你生活在843年的凡尔登，你会怎么看待三兄弟分家的决定？
  \item 如果欧洲也像秦朝一样统一了文字，历史会不一样吗？为什么一个统一的政治实体几乎一定需要一个统一的书写系统？
\end{enumerate}

\vspace{0.3cm}
\textcolor{blue!70!black}{\textbf{延伸阅读}}
\begin{itemize}
  \item 可以回看《中国古代史》第一课的"中西封建对比"和第三课的"推恩令"——不同的继承制度，不同的历史后果
\end{itemize}

\newpage

\section{\textcolor{orange!80!black}{第二课\quad 皇帝 vs 教皇：谁能任命主教？}}

\subsection*{一、地下的博弈：谁控制教会？}

在上一课我们看到，查理曼帝国分裂之后，西欧变成了大大小小的封建领地。但在这些领主的头上，还存在一个跨越国界的组织——\textbf{基督教会}。谁控制教会，谁就控制了一个遍布全欧洲的\:\textbf{"精神行政网络"}。教区的神父能让农民服从领主，也能让他们反抗——取决于神父听谁的。

这个地方的主教应该由谁来任命？这是个看似无聊的宗教细节，实际上是一个让中世纪的皇帝和教皇撕破脸、打了半个世纪的核心权力问题。

\subsection*{二、亨利四世跪在雪地里：一场没有赢家的博弈}

1075年，教皇\textbf{格里高利七世}发布了一道命令：\textbf{以后只有教皇有权任命主教，世俗统治者——包括皇帝——不得干预}。这是一个炸弹。因为神圣罗马帝国的皇帝长期以来都把主教职位当作政治工具——把亲戚和心腹安插在富庶的教区，既能控制财源，又能巩固统治。

时年25岁的神圣罗马帝国皇帝\textbf{亨利四世}收到这封信之后怒火中烧。他回了一封措辞极度侮辱的信，开头是"\textbf{亨利，非由于篡夺而是由于上帝恩典的国王，致希尔德布兰德——不是教皇，而是一个假修士……}"。他召集了忠于他的德意志主教开会，宣布罢免格里高利七世。

格里高利七世的回应堪称中世纪政治史上最狠的一招：他宣布将亨利四世\textbf{开除教籍}。在一个人们深信开除教籍的国王死后一定会下地狱的时代，这条命令的影响力是现代人难以想象的。亨利的封臣们\u2014\u2014本来就不满这个年轻高傲的国王\u2014\u2014找到了一个完美的借口：我们不能服从一个被开除教籍的人。大批贵族叛变。

亨利发现如果他不低头，他的王国会在几周之内解体。1077年1月，他带着妻子和少量随从翻越了阿尔卑斯山（那年的冬天格外寒冷），到达教皇当时所在的\textbf{卡诺莎城堡}。教皇拒绝见他。亨利在城堡门外站了\u2014\u2014根据史料的记载\u2014\u2014\:\textbf{三天}。据说他赤着脚，穿着悔罪者才穿的粗毛衣服，在隆冬的寒风中站在雪地里。

最终格里高利七世接见了他，恢复了教籍。\textbf{卡诺莎之辱}——一个皇帝跪在教皇脚下的画面\u2014\u2014从此成为欧洲政治史上的一个标志性场景。直到今天，德语里仍然有"走一趟卡诺莎"这种说法（nach Canossa gehen），用来形容被迫低头认错。

\subsection*{三、斗争的结果：谁也没赢透}

卡诺莎的戏剧性很容易让人误以为"教皇赢了"。实际上根本不是。

亨利四世回国之后迅速收拾了叛变的贵族，转过头来就和教皇再度开战。他带兵攻入了意大利，赶走了格里高利七世。1085年，格里高利在流亡中死去，他的遗言被后人记录为："\textbf{我热爱正义，憎恨邪恶——所以我死于流放。}"

教皇和皇帝的这场争斗持续了将近半个世纪，最终在1122年达成了\textbf{沃尔姆斯宗教协定}。这份协定本质上是一个各退一步的妥协：皇帝不再任命主教，但主教在就任教职之前必须先向皇帝宣誓效忠。双方谁都没有得到自己想要的全部东西。

\shihai{拜占庭："政教合一"的另一条路}{
  正在西欧的皇帝和教皇打得不可开交的时候，东边的拜占庭帝国（东罗马）走了一条完全不同的路。

  在拜占庭，\textbf{皇帝是教会的最高领袖}。君士坦丁堡的牧首（相当于东正教的教皇）由皇帝任命，不能独立于皇权。这套制度叫做"\textbf{政教合一}"——和西欧正在形成的"政教分离"形成鲜明对比。

  这两种模式哪一种更好？答案取决于你问谁。政教合一的好处是稳定——拜占庭帝国存活了一千年，期间极少出现西欧那种教皇废黜皇帝的极端案例。坏处是\u2014\u2014当统治者的决策和教会的信仰相悖时，没有独立的教权来形成制衡。而西欧的政教分离模式带来了持续的内部冲突（皇帝和教皇的争斗只是一个开始），但它在\textbf{意外中}创造了一个重要的副产品：\textbf{没有任何一个单一的权力中心可以垄断一切}。当没有人是绝对的权威时，不同的势力就被迫互相谈判、签协定、建立规则——这些"被迫的协议"构成了后来西欧宪政和法律传统的早期渊源。

  再次强调：这不是因为西欧的皇帝和教皇比拜占庭人更明智。只是拜占庭的皇帝一开始就足够强大，吞掉了教权；而西欧的皇帝一开始就不够强大，吞不掉——所以他们只好互相妥协。政治制度的差异，很多时候不过是\textbf{力量对比}的副产品。
}

\subsection*{本课要点}

\begin{center}
\begin{tabular}{|c|c|p{7cm}|}
\hline
\textbf{事件} & \textbf{年代} & \textbf{意义} \\
\hline
格里高利七世教皇令 & 1075年 & 宣布只有教皇有权任命主教 \\
\hline
卡诺莎之辱 & 1077年 & 亨利四世在教皇门外跪了三天 \\
\hline
沃尔姆斯宗教协定 & 1122年 & 各退一步：主教由教会选任，但须向皇帝效忠 \\
\hline
\end{tabular}
\end{center}

\vspace{0.3cm}
{\large\textcolor{orange!80!black}{\textbf{想一想}}}
\begin{enumerate}
  \item 格里高利七世不是靠武力而是靠"开除教籍"逼迫了皇帝低头——这种"软的权力"为什么在11世纪如此有效？如果今天有一个强大的组织声称谁不听话就开除他，还有人会在乎吗？
  \item 西欧和拜占庭走了两条不同的路——政教分离和政教合一。读完之后，你认为哪种模式更稳定？更"自由"的一定更好？
\end{enumerate}

\newpage

\section{\textcolor{orange!80!black}{第三课\quad 西班牙：八百年的收复失地}}

\subsection*{一、七个人退回山洞}

711年，一支来自北非的穆斯林军队渡过直布罗陀海峡，在短短几年之内征服了几乎整个伊比利亚半岛。西哥特王国——此前统治西班牙的日耳曼政权——在初次交战中就被击溃了。穆斯林征服者建立的政权在西班牙南部的科尔多瓦建都，此地后来发展成为当时西欧最富裕、最文明的城市（同期的巴黎和伦敦相比之下只是大型村庄）。

但在北部的阿斯图里亚斯山区，有七个西哥特贵族拒绝投降。他们退入了坎塔布连山脉，在科瓦东加的一个山洞里做了一个被后来的西班牙民族主义叙事极度神圣化的动作——他们打了一场小规模的伏击战，杀死了当地的一支穆斯林驻军，宣布了基督教西班牙的存续。从718年科瓦东加之战算起，到1492年最后一座穆斯林据点格拉纳达沦陷——\textbf{收复失地运动}（Reconquista）整整进行了将近八百年。

\subsection*{二、八百年的边疆社会}

收复失地不是一个持续不断的"战争"，而是一个断断续续的过程。基督教王国和穆斯林政权之间经常有长达几十年的和平期，彼此通商、通婚、交换文化。但有一个结构性的事实从未改变：伊比利亚半岛的基督教王国\textbf{活在一个边疆地带}——往南一步就是"敌方"的地盘。这种长期的、延续了数百年的边疆状态，对西班牙社会产生了深刻的影响。

在"收复失地"的过程中，国王需要贵族和骑士冲锋陷阵，贵族需要国王授予他们从穆斯林手中夺来的新土地。因此，和同时期的法国比起来，西班牙的国王在\textbf{早期就比较强大}——他没有被一群拥有独立领地的傲慢大贵族架空，因为他不断地有新土地可以赏赐那些为他战斗的人。

而收复的土地上，有一个非常现实的治理问题：\textbf{刚刚拿下来的城镇，里面住着穆斯林和犹太人，谁来管？}西班牙国王的答案是教会。把镇子交给教会——修道士经营土地、神父管理社区、大主教为国王的远征提供精神合法性。教会和国王形成了一种\textbf{牢固的同盟}：国王给教会土地和权力，教会给国王统治的合法性。这解释了为什么后来的西班牙成为了全欧洲最狂热的天主教国家。

\subsection*{三、1492：同一个年份里的三件事}

1492年对于西班牙来说是一个极端浓缩的历史转折点。一月，穆斯林在伊比利亚半岛上最后的堡垒——格拉纳达——向\textbf{费尔南多}和\textbf{伊莎贝拉}投降。这对夫妻分别是阿拉贡和卡斯蒂利亚的君主，他们的婚姻使西班牙实质上完成了统一。收复失地运动在八百年之后正式宣告结束。

三月，费尔南多和伊莎贝拉签署了《阿尔罕布拉法令》——所有拒绝改信基督教的犹太人必须在四个月内离开西班牙。据估计，大约有十万到二十万犹太人被驱逐——其中不乏医生、学者、银行家和大商人。西班牙在宗教上"清洗"了自己的社会，代价是流失了大量的经济和智力资本。

八月，热那亚航海家\textbf{克里斯托弗·哥伦布}——在费尔南多和伊莎贝拉的资助下——率领三艘小船从西班牙西南部的帕洛斯港出发，向西横渡大西洋。十月十二日，他发现了陆地（他以为到了亚洲的印度群岛，实际上到了今天的巴哈马）。

这三件事情不是时间上的巧合。收复失地运动塑造了一个什么样的西班牙？一个以天主教为精神核心、以国王为政治中心、以战争和征服为国家经验的西班牙。统一之后，国王有了一支经验丰富的常备军没事干，有一批在边疆战事中成长起来的贵族渴望新的军功，有一个刚刚被"净化"的宗教狂热想要向外输出。\textbf{把这种能量导向大西洋}——几乎是顺理成章的。

\subsection*{四、本课要点}

\begin{center}
\begin{tabular}{|c|c|p{7cm}|}
\hline
\textbf{事件} & \textbf{年代} & \textbf{意义} \\
\hline
穆斯林征服伊比利亚 & 711年 & 西哥特王国崩溃；科尔多瓦成为西欧最璀璨的城市 \\
\hline
收复失地运动 & 718—1492年 & 持续近800年的边疆战争；塑造了西班牙的军事贵族文化和政教同盟 \\
\hline
格拉纳达陷落 & 1492年1月 & 收复失地终结；西班牙统一 \\
\hline
驱逐犹太人 & 1492年3月 & 约10—20万犹太人被逐；经济和文化人才的重大损失 \\
\hline
哥伦布远航 & 1492年8月 & 同一个国家，同一年份，把内陆战争的能量转向了海洋 \\
\hline
\end{tabular}
\end{center}

\vspace{0.3cm}
{\large\textcolor{orange!80!black}{\textbf{想一想}}}
\begin{enumerate}
  \item 收复失地运动持续了八百年——比从宋朝到今天的时间还要长。一个社会持续处于"边疆战争"状态，会对它的制度、文化和民族性格产生什么影响？
  \item 1492年同一年里，西班牙驱逐了犹太人，又资助了哥伦布——驱逐和远航是一体两面吗？
\end{enumerate}

\vspace{0.3cm}
\textcolor{blue!70!black}{\textbf{延伸阅读}}
\begin{itemize}
  \item 纪录片《西班牙：从收复失地到帝国》——BBC出品
  \item 如果想了解1492年前后更宏大的全球图景，可以回看《世界古近代史》第五课"大航海时代"
\end{itemize}

\newpage

\section{\textcolor{orange!80!black}{第四课\quad 英格兰：法律如何一步一步管住了国王}}

\subsection*{一、诺曼人打下了一个效率惊人的封建机器}

1066年是英格兰历史上最彻底的"系统重装"年份。这一年，诺曼底公爵威廉率领一支说古法语的军队渡过英吉利海峡，在哈斯廷斯战役中击败了英格兰国王哈罗德的军队，哈罗德本人被一箭射穿眼睛阵亡。威廉成为了英格兰的国王，史称"\textbf{征服者威廉}"。

威廉做了一件此前任何一个英格兰国王都没能做到的事：他派人走遍了全国，把每一块土地、每一个庄园、每一个磨坊、每一个农民——甚至每一头猪——都登记在了一本卷宗里。这就是著名的《\textbf{末日审判书}》。这本册子的目的非常简单：\textbf{征税}。国王现在知道了这个国家有多少财产，该向谁收税、收多少。在此之前的封建欧洲，国王是靠封臣的"效忠"和"兵役"来维持统治的——你的兵是你的兵，你的封臣的兵是封臣的兵。威廉把这套系统改造成了一个以国王为顶点的\textbf{垂直管理体系}。

在法国，国王的土地比好几个大贵族还小，国王的权威常常出不了巴黎。而在英格兰，征服者威廉直接把全国的封建体系搭建成了一张以国王为核心的放射状网络——\textbf{英格兰从"零"开始的时候就被建成了一个中央集权程度超过同时期任何西欧国家的封建机器。}

\subsection*{二、约翰王：一个糟糕的国王碰上了一个好时机}

这个高效的集权机器在运转了大约一百五十年之后，遇上了一个致命的麻烦：\textbf{国王太烂了}。

\textbf{约翰王}——也就是罗宾汉传说里那个贪婪无能的"坏国王"的原型——在1199年登上王位之后做了一系列教科书级的错误：对法战争连连失利，把英格兰在法国的大片领地丢了个干净（从此他得了个外号叫"无地者约翰"——用国王丢掉国土来命名，羞辱性极强）；为了打仗疯狂加税，贵族们被榨得喘不过气；和教皇英诺森三世闹翻，被开除教籍，搞到整个英格兰被置于"\ruby{禁}{jìn}\ruby{令}{lìng}"之下——整整六年教堂不开放、不举办弥撒、不为死者举行葬礼。

贵族们忍无可忍了。1215年6月15日，一批全副武装的男爵在兰尼米德草地上逼着约翰王在一份文件上盖了章。这份文件叫《\textbf{大宪章}》。

\subsection*{三、大宪章不是什么？}

大宪章在后来的历史教科书里经常被描写为"\textbf{自由的开端}"或者"英国民主的基石"。但如果你知道这些逼着国王签字的贵族是什么人，你就会重新看待这段话。

1215年的大宪章\textbf{和普通农民没有一丁点关系}。它是盎格鲁—诺曼大贵族为了保护自己的\textbf{既有特权}而逼国王签字的一份法律文件。它的核心条款——"国王未经王国共同协商不得征收额外赋税"——字面上看起来冠冕堂皇，实际上翻译过来是：\textbf{国王你以后想加我们贵族的税，必须经过我们贵族自己开会同意}。

但是这个文件里确实出现了一些在法律史上具有深远影响的表述。比如第三十九条——"\textbf{任何自由人未经同等地位者依法审判或国家法律裁决，不得被逮捕、监禁、剥夺财产、流放或以任何方式毁灭。}"这就是后世"\textbf{正当法律程序}"的胚胎——虽然这里的"自由人"只占当时英格兰人口的一小部分（大量的农奴不在此列），但它确实是对"\textbf{国王不能仅凭个人好恶就抓人}"这条原则最早的成文法表达之一。

\subsection*{四、砍了国王的头，又把国王请回来}

大宪章只是开端。此后几百年里，英国国王和议会之间经历了一个漫长的\textbf{反复拉锯}过程，没有任何人可以提前计算出结局。

\begin{itemize}
  \item \textbf{13世纪后期，议会成型。} 爱德华一世为了筹措对苏格兰和威尔士的战争费用，频繁召集贵族、教士和平民代表开会讨论征税——这个为了\textbf{搞钱}才召集的会议，慢慢演化成了\textbf{议会}。国王需要钱$\to$议会要求权利$\to$国王答应条件$\to$议会批准税收——这是一个\textbf{纯粹的利益交换过程}，不是某个政治哲学家的理想设计。

  \item \textbf{17世纪，矛盾爆发。} 查理一世试图抛开议会单独统治，1629—1640年间他根本不召开议会（"\textbf{十一年暴政}"）。1642年，议会的武装和国王的武装终于打了起来——英国内战爆发。1649年，议会军在奥利弗·克伦威尔的领导下获胜。查理一世被公开审判并斩首——\textbf{一个国王被自己的臣民以法律程序处死}，这件事震撼了整个欧洲的君主们。英国宣布成为"\textbf{共和国}"。

  \item \textbf{然后又请回了国王。} 克伦威尔死后，共和国的军事统治迅速失控，没有人能替代他。1660年，英格兰的议会\textbf{主动邀请}查理一世的儿子查理二世从流亡地法国回来做国王——\textbf{反对专制不等于反对君主制}。英国人不想要一个不受限制的国王，不意味着他们想要一个没有国王的军事独裁。

  \item \textbf{1688年，再赶走一个换一个。} 查理二世的弟弟詹姆斯二世公开信奉天主教——在一个新教已经成为国教的国家，一个天主教国王让所有人都睡不着觉。议会秘密邀请詹姆斯二世的女儿玛丽和她的丈夫——荷兰的奥兰治亲王威廉——来"接管"英格兰。威廉带兵登陆，詹姆斯逃往法国。议会宣布他"放弃了王位"，请威廉和玛丽共同即位。这就是被后世称为"\textbf{光荣革命}"的事件。
\end{itemize}

\subsection*{五、"光荣革命"之后——并没有那么"光荣"}

光荣革命在维多利亚时代的辉格派历史学家笔下被描写为"英国人的自由从此永固"——这是一个有问题的叙事。

1689年议会通过了《权利法案》，国王未经议会同意不得征税、不得在和平时期维持常备军、议员在议会中有言论自由（"议会中的言论自由、辩论或议事之自由，不应在议会之外的任何法院或任何地方受到弹\rube{劾}{hé}或讯问"）。这些条款确实约束了王权。

但\textbf{约束≠王权就成了摆设}。18世纪的国王——特别是乔治一世、乔治二世和乔治三世——仍然通过任命首相、收买议员、操控选举等各种手段实质性地干预政治。这一时期首相多由国王信任的人担任（沃波尔长期掌权很大程度上就是因为乔治一世不懂英语、懒得参与内阁会议，给了沃波尔自行其是的机会）。直到19世纪，随着选举权的逐步扩大（1832年改革法案、1867年改革法案），随着政党内阁制的成熟，国王才真正退到了我们现代人熟悉的"虚君"位置上。

\textbf{英国的君主立宪不是一个"设计好"的制度，是一层一层叠出来的惯例}——大宪章、内战、共和国失败、复辟、光荣革命、选举改革——没有一个时间点是"完成态"。每一步都是前一步出了问题之后的修补。

\shihai{波兰的"自由否决权"——制约权力走到极致就是亡国}{
  就在英国一步步探索怎么约束王权的时候，东边的波兰—立陶宛联邦也在进行一场\textbf{截然不同的"权力制约"实验}。

  波兰的贵族（施拉赤塔）从16世纪开始获得了对王权极其强大的制约能力——其中最出名（也最致命）的是"自由否决权"（liberum veto）。1652年，波兰议会首次行使了这一权力：任何一个议员，只要站起来喊一声"我反对"，就能让整个议会的决议作废——他不需要说服别人，不需要任何理由。

  这项制度的初衷是保护少数派贵族的权利不被国王或多数派践踏——听起来很高尚。但实际的结果是\textbf{国家瘫痪了}。几个世纪以来，任何一个外国势力只要收买一个波兰议员，他走进议会说一句"我反对"，整个国家的决策就停摆了。从1652年到1764年，波兰议会召开了71次，有42次——超过一半——都因为这一个否决权而无法做出任何决定。

  1772年、1793年、1795年，俄国、普鲁士和奥地利\textbf{三次瓜分}了波兰。1795年波兰从地图上彻底消失了，直到1918年才重新出现。一个制度设计用来的防止国王专制的工具，最终让整个国家连自己都保卫不了。

  \textbf{英国和波兰走了两条几乎相反的路。}英国面对的问题是"怎样让国王不能为所欲为"——它的答案是\textbf{逐渐建立了一个有实权的议会}，用集体决策代替个人专制。波兰面对同样的问题——它的答案是\textbf{每个贵族都可以独自否决所有人的决定}。结果：英国在18—19世纪崛起为全球帝国，波兰连自己的领土都没能守住。\textbf{约束权力不等于让权力瘫痪——这两者之间的界限，一不小心就会踩过了线。}
}

\subsection*{六、本课要点}

\begin{center}
\begin{tabular}{|c|c|p{7cm}|}
\hline
\textbf{事件} & \textbf{年代} & \textbf{关键信息} \\
\hline
诺曼征服 & 1066年 & 威廉建立高效集权的封建体系；《末日审判书》 \\
\hline
大宪章 & 1215年 & 贵族逼国王签字保护自己特权；"正当法律程序"胚胎 \\
\hline
英国内战+查理一世被斩首 & 1642—1649年 & 国王与议会的武装冲突；斩杀国王震惊欧洲 \\
\hline
王朝复辟 & 1660年 & 共和国失败；迎回查理二世 \\
\hline
光荣革命 & 1688年 & 再换一个国王；权利法案（1689） \\
\hline
君主立宪制的真正成熟 & 18—19世纪 & 选举权扩大、责任内阁制形成——不是一次"革命"完成的 \\
\hline
\end{tabular}
\end{center}

\vspace{0.3cm}
{\large\textcolor{orange!80!black}{\textbf{想一想}}}
\begin{enumerate}
  \item 英国从大宪章到权利法案经历了四百多年——约束国王的权力为什么需要这么久？如果一步到位——"明天起国王不许做决策"——会怎样？
  \item 英国约束了王权却没出事，波兰约束了王权却亡了国——同样都是在"限制权力"，区别在哪里？
\end{enumerate}

\vspace{0.3cm}
\textcolor{blue!70!black}{\textbf{延伸阅读}}
\begin{itemize}
  \item 《大宪章》的现代英文译本（大英图书馆官网可以免费阅读）
  \item 纪录片《君主制：英国王室史》——BBC出品，追踪英国王权演变的完整历程
\end{itemize}

\newpage

\section{\textcolor{orange!80!black}{第五课\quad 百年战争：英法恩怨与民族国家的诞生}}

\subsection*{一、诺曼底的公爵怎么成了英国的国王？}

我们在上一课英国篇里讲过1066年诺曼底公爵威廉征服英格兰的故事。这个故事有一个极其别扭的后果，让英法两国彼此折磨了四个世纪：\textbf{英格兰的国王同时也是法国的封臣。}威廉征服英国之后，保留了诺曼底公爵的封号和领地，向法国国王行封臣礼。这意味着一个奇怪的场景：在英格兰他是至高无上的国王，在诺曼底他却是法国国王的"下属"——名义上要向法王效忠。

更麻烦的是，通过联姻和继承，英国王室在法国控制的地盘越来越大。到12世纪中叶，金雀花王朝的亨利二世除了是英格兰国王之外，还通过联姻控制着诺曼底、安茹、阿基坦等大片法国西部土地——英国国王在法国的领地，比法国国王自己直接控制的土地大好几倍。这就好比你是公司的董事长，但你手下某个"部门经理"实际上去年兼并的产业比你自己的还多。

法国国王不能忍。到了14世纪，法国卡佩王朝的主支绝嗣，英国国王爱德华三世（他是前法王腓力四世的外孙）宣布自己对法国王位有继承权。这等于说"\textbf{我不光要你那大片地，我要你的整个王国}"。

\subsection*{二、打了不止一百年}

1337年，百年战争正式爆发。但是说"打了一百多年"容易误导——这不是一场连续不断的战争。它更像是一百多年里断断续续爆发的一系列战役、围城和劫掠，中间夹杂着很多次的停战和瘟疫（黑死病——当时人死了一大批，谁还有精力打仗）。

战况极度不均衡。战争的前八十年里，英国人几乎场场获胜——\textbf{克雷西战役}（1346年）和\textbf{阿金库尔战役}（1415年）都是军事史上以少胜多的经典案例。英国为什么这么能打？一个重要原因是英国军队的结构和法国完全不同。英国军队的核心是\textbf{长弓手}——一群训练有素的职业步兵，能在一分钟内射出十到十二支穿甲箭。法国军队的核心是重装骑士——贵族们骑着高头大马、穿着重甲冲锋。但当长弓手的箭雨落下时，马匹被射翻、骑士倒地后被自己的盔甲困住——法国的贵族阶层在克雷西和阿金库尔遭到了惨重的损失。

法国这边就更惨了：前线不断输，后方四分五裂——勃艮第公爵公开和英国人结盟，法国王室自己的亲戚互相背叛——\textbf{法国打了八十年，快把自己打成碎片了}。

\subsection*{三、一个17岁女孩的登场}

1429年法国已经到了悬崖边上。英军围困了卢瓦尔河畔的奥尔良城——这座城一旦失守，整个法国南部将门户大开。法国王太子查理连正式的加冕礼都办不了（因为加冕地兰斯被英军控制），他被讥讽为"\textbf{布尔日的土皇帝}"——地盘只有巴掌大小。

就在这个节骨眼上，一个17岁的农村女孩走进了查理的宫廷。她告诉王太子：她听到了上帝的声音，她被选中来拯救法国，她要带领军队解除奥尔良的围困，护送他去兰斯加冕。

让人意外的是——绝望中的查理同意给她盔甲、武器和一支军队。

关于\textbf{贞德}（Jeanne d'Arc），我们分析两个关键问题。第一，贞德实际上\textbf{几乎不直接参与战斗}——她主要在战场上举着她的军旗（上面写着"耶稣马利亚"，而非法兰西），站在前线，让所有士兵看到她。她告诉进攻的士兵"往前走，上帝和我们在一起"——而在此之前，奥尔良的守军在城里足足缩了七个月没敢出城。第二，贞德对逃兵的零容忍——她看到了溃退的法军就策马追上去、用剑背抽他们逼回去——"有上帝帮助你们，往回跑什么？！"贞德本质上是担任了一个\textbf{心理学上的引爆点}：在一个士气已经低到地底的军队里，出现了一个声称和上帝有直接热线的小姑娘。这比任何军事策略都更能鼓舞一支绝望的军队。

1429年5月，奥尔良围城被解除——在贞德到来之后的九天之内。整个法国战场上的士气发生了根本性的逆转。一个月内，法军一路向前推进。7月17日，贞德护送查理七世抵达兰斯，完成了加冕。

\subsection*{四、她被烧死的原因}

1430年5月，贞德在一次小规模战斗中被俘。把她俘虏的勃艮第人显然认为贞德"值大钱"——把她卖给了英国人。英国人花钱买她只有一个目的：\textbf{让她从"上帝的人"变成"女巫"}。如果贞德真的是神在指引她，那上帝就站在法国一边——英国人从宗教逻辑上就输了。他们必须证明贞德的"声音"不是来自上帝，而是来自魔鬼。

鲁昂的宗教审判（1431年）——一群受过顶级教育的巴黎大学神学教授审问一个不会读书的19岁农家女。庭审记录保留了贞德的很多回答。当被问及"你在上帝面前是否感到自己是恩典状态？"时——这是一个神学界的陷阱问题。如果你说是（而教会的判断认为你不是），你就说了假话；如果你说不确定，你就等于承认了自己可能犯了罪。贞德的回答堪称经典：\textbf{"如果我不是，愿上帝把我放在那里；如果我是在恩典中，愿上帝把我留在此中。如果我是，而我说不是，我就是撒了谎；如果我不是，而我说我是，我就是骄傲。"}在场的教士当场愣住了：一个不识字的乡下姑娘，给出了一个连神学博士都挑不出毛病的回答。

但是这个审判的结论是早就预定好的。1431年5月30日，贞德在鲁昂老市场广场被绑在火刑柱上烧死。她在一个异端帽上系了一张纸条，上面只是三个词：\textbf{Jhesus Maria}——耶稣，马利亚。

1453年，贞德死后22年，百年战争以法国的胜利告终。英国在大陆上的领地几乎全部丧失，只保留了加莱——一个小小的港口，一百多年后也丢了。

\subsection*{五、百年战争的意外后果}

百年战争最重要的遗产可能不是领土的得失，而是\textbf{人们开始觉得自己是"法国人"或"英国人"}。

在战争之前，中世纪欧洲的普通农民对自己的身份认同是很模糊的：你是"属于某个领主封地的人"，而很难说是"法国人"——毕竟你的领主可能是一个说法语的英国贵族，或者一个效忠法王的德意志骑士。但当战争以"法国vs英国"的形式打了整整一个世纪——英国的长弓手在法国乡间烧杀抢掠、法国的士兵在英国的属地征收军饷——普通人的认知开始悄悄改变：\textbf{"他们"和"我们"不一样。}

法国在这场战争中获得了一个"共同的敌人"。面对英国入侵和内战的双重压力，法国各个阶层被迫整合、统一在一个国王之下。到15世纪末，法国从一个四分五裂的封建领地集合，进化成了一个早期\textbf{民族国家}——这是中世纪西欧最早、最有凝聚力的统一君主国之一。

另外军事上，火药武器在百年战争的后期开始出现在战场上（卡拉维勒战役中使用了早期火炮）。火药爆炸的巨响和强大的威力，比长弓更彻底地宣告了城堡和骑士时代的终结——躲在石墙后的领主不再安全，国王可以用火炮轰开任何拒绝交税的贵族堡垒。

\subsection*{六、本课要点}

\begin{center}
\begin{tabular}{|c|c|p{7cm}|}
\hline
\textbf{事件} & \textbf{年代} & \textbf{关键信息} \\
\hline
百年战争爆发 & 1337年 & 英国国王索要法国王位；不仅是王位之争，更是法国领土控制权之争 \\
\hline
克雷西/阿金库尔 & 1346年/1415年 & 英军长弓以少胜多；法国贵族阶层受重创 \\
\hline
贞德解奥尔良之围 & 1429年5月 & 九天内逆转整个战争态势 \\
\hline
贞德被处死 & 1431年5月 & 英方以宗教审判"证明"她是女巫 \\
\hline
战争结束 & 1453年 & 英国丧失除加莱外的全部大陆领地 \\
\hline
民族认同成型 & 15世纪末 & 战争催生了法国和英国的最初民族意识 \\
\hline
\end{tabular}
\end{center}

\vspace{0.3cm}
{\large\textcolor{orange!80!black}{\textbf{想一想}}}
\begin{enumerate}
  \item 贞德不懂军事策略却扭转了战局——她改变了什么？一支军队里士气的作用到底有多大？
  \item 百年战争之前，一个法国南部农民可能从来没想过自己是"法国人"。战争如何改变了这种状态？"我们"的定义是由什么建立起来的？
\end{enumerate}

\vspace{0.3cm}
\textcolor{blue!70!black}{\textbf{延伸阅读}}
\begin{itemize}
  \item 莎士比亚戏剧《亨利五世》——虽然文学化了阿金库尔战役，但"我们少数，我们快乐的少数"这段演讲已经成为英语文学中最经典的战前动员
  \item 《贞德传》少儿版——无论如何，这个故事本身比任何小说都更有魔力
\end{itemize}

\newpage

\section{\textcolor{orange!80!black}{第六课\quad 三十年战争：主权国家的规则是怎样定下来的}}

\subsection*{一、布拉格的"扔出窗外"}

1618年5月23日，布拉格。一群波希米亚（今捷克）的新教贵族怒气冲冲地冲进布拉格城堡，找到了两位哈布斯堡皇帝派来的总督和他们的秘书。这几位来客是来推行反新教政策的——波希米亚的新教徒对此忍耐到了极点。贵族们用简单直接的方式表达了自己的态度：他们把三位官员从城堡窗户扔了出去。

这三个倒霉的人从十几米高处摔进了城堡的壕沟里。但他们没死——掉在了壕沟里堆着的垃圾堆上（一说是马粪堆）。天主教方面后来声称他们是被天使托住了，新教方面则说他们是掉进粪堆捡了条命。

不管怎么说，\textbf{布拉格掷出窗外事件}引爆了一场战争——\textbf{三十年战争}（1618—1648年）。它是欧洲历史上破坏性最大的战争之一，持续整整三十年，把德意志地区打成了地狱。

\subsection*{二、从宗教战争到霸权战争}

三十年战争的一个核心特征是你很难用一句话说清楚"\textbf{到底是谁在打谁}"——因为参战方随着时间在不断地换阵营。

最初是波希米亚新教徒反哈布斯堡天主教皇帝——这是宗教战争。但后面的事情迅速复杂化：信天主教的法国加入了\textbf{新教阵营}，信新教的丹麦和瑞典互相争抢北德地盘，西班牙（哈布斯堡家族的另一分支）出兵支持奥地利哈布斯堡。\textbf{宗教的外衣裹不住赤裸裸的利益计算}。

德意志地区在这场战争中被反复碾压。三十年里，雇佣军来来回回地扫荡同一片地区——你驻扎在农村，农民就得供你吃住；你走了，另一支军队又来了。据估计，\textbf{德意志地区在这场战争中损失了约三分之一到二分之一的人口}——部分地区直接变成了无人区。这个人口损失比例，和黑死病相当。战争中被焚毁的村庄、城镇、教堂数不胜数——德意志的恢复时间长达一个世纪以上。

\subsection*{三、红衣主教和"国家利益"}

要理解三十年战争是怎么结束的，需要把目光从战场上转向外交桌上。这里有一个人不得不提：法国的\textbf{红衣主教黎塞留}。

黎塞留是天主教会的枢机主教——理论上他应该不遗余力地支持天主教阵营。但是他脑子里装的不是"宗教信仰"，而是\textbf{法国的国家利益}。从法国的角度看，最可怕的敌人不是远方的新教国家（丹麦、瑞典），而是把法国包围起来的哈布斯堡家族——这个家族同时控制着西班牙和神圣罗马帝国，把法国东西两边都堵住了。\textbf{一个统一的、由哈布斯堡家族控制的天主教欧洲，对法国来说是一个灾难}——不管它有多么"天主教"。所以黎塞留做了一个让当时的欧洲人目瞪口呆的决定：\textbf{法国出钱支持新教瑞典的军队去打天主教的哈布斯堡皇帝}。

黎塞留自己没活到战争结束——他1642年去世，三十九天后法王路易十三也死了。四年后战争结束，法国的位置稳如磐石——哈布斯堡的包围圈被打碎了。

\shihai{威斯特伐利亚和约——现代国际关系的操作手册}{
  1648年，各方终于坐下来签订了《\textbf{威斯特伐利亚和约}》。这不是一份"和平协议"那么简单——它确立了一套全新的国际关系游戏规则，至今仍在世界政治中运作。

  这套规则的几个核心原则是：
  \begin{itemize}
    \item \textbf{主权平等}：每一个签署国（不管大小）在自己的领土内享有最高权力。其他国家没有权力因为宗教、意识形态或者"某种更崇高的理由"来干涉一个主权国家的内政。
    \item \textbf{不干涉内政}：一个国家怎么治理——信仰什么宗教、用什么法律、国王是谁——是自己的事。\textbf{其他主权国家无权以任何理由强行干涉}。
  \end{itemize}

  你可能会觉得，这些规则今天不都是常识吗？在1648年之前，这种规则是不存在的。在中世纪，教皇\textbf{确实}会召集军队去攻打一个"不听话"的基督教国王；在国际关系里，信仰和"更高的道德理由"\textbf{确实是}干涉别国的合法理由（十字军东征就是这种逻辑下的产物）。威斯特伐利亚体系把这种逻辑\textbf{终结了}：你说你是为了"上帝"还是为了"正义"——拿条约来，拿公认的边界来。

  今天联合国宪章里"不干涉内政"的原则，可以一路追溯到1648年的这份和约。而三十年战争告诉人们的是：打出这个原则的过程——付出了整个德意志三分之一到一半人口的代价。\textbf{现代世界里的每一条基本规则，背后几乎都埋着战争的白骨。}
}

\subsection*{四、神圣罗马帝国的实质终结}

三十年战争还给欧洲政治版图带来了一个重大变化：\textbf{神圣罗马帝国从一个"帝国"退化成了一个地理名词}。

和约规定，帝国内的每一个邦国——普鲁士、巴伐利亚、萨克森、勃兰登堡……——都有权利独立地和外国签订条约、宣战、建立军队。这些"邦国"实际上变成了独立的主权国家，帝国的皇帝只剩下了名义上的尊号。此后两百多年里，神圣罗马帝国一直是一个"叫做帝国但不像帝国"的空壳——直到1806年被拿破仑正式解散。

中欧从此由一个大帝国变成了几百个大小不等的独立邦国——这种分裂状态一直延续到19世纪后半叶，普鲁士通过一系列的战争（1864年对丹麦、1866年对奥地利、1870—1871年对法国）才把德意志重新统一起来——但那是另一本书的内容了。

\subsection*{五、本课要点}

\begin{center}
\begin{tabular}{|c|c|p{7cm}|}
\hline
\textbf{事件} & \textbf{年代} & \textbf{要点} \\
\hline
布拉格掷出窗外 & 1618年 & 波希米亚新教徒反哈布斯堡；引爆三十年战争 \\
\hline
战争全面化 & 1620s—1630s & 宗教冲突转变为全欧争霸；天主教法国支持新教瑞典 \\
\hline
威斯特伐利亚和约 & 1648年 & 确立主权平等 + 不干涉内政原则；神圣罗马帝国名存实亡 \\
\hline
\end{tabular}
\end{center}

\vspace{0.3cm}
{\large\textcolor{orange!80!black}{\textbf{想一想}}}
\begin{enumerate}
  \item 黎塞留是天主教的红衣主教，却出钱支持新教军队打天主教的皇帝——他做错了还是做对了？如果在"信仰"和"国家利益"之间必须做出选择，怎么选才合理？
  \item 威斯特伐利亚和约里的"不干涉内政"原则，到今天还在被各国反复引用。一个三百多年前定下来的规则在全球化的今天还适用吗？还是需要被修改？
\end{enumerate}

\vspace{0.3cm}
\textcolor{blue!70!black}{\textbf{延伸阅读}}
\begin{itemize}
  \item 纪录片《三十年战争：铁与火的时代》——用视觉效果还原这场灾难
  \item 彼得·威尔逊《欧洲的心脏：神圣罗马帝国史》（有通俗版）——理解这个奇怪的政治体为什么能"活着"这么久
\end{itemize}

\newpage

\section{\textcolor{orange!80!black}{第七课\quad 理性革命：从听上帝到信自己}}

\subsection*{一、一颗行星改变了一切}

1543年，一个波兰教士\textbf{尼古拉·哥白尼}在临终的病床上拿到了一本刚刚印出来的书——《天体运行论》。他用拉丁文就着复杂的数学推导论证了一个在今天看来毫无争议、但在当时会把人送上火刑柱的结论：\textbf{不是太阳围着地球转，而是地球围着太阳转}。

哥白尼本人是天主教会的教士——他从来没有"反对宗教"。他把这本书题献给了教皇。但出版人奥西安德可能是怕惹麻烦，在书的序言里擅自加了一段话："不要把这些计算当成物理现实——它们只是一个有用的数学模型。"哥白尼死前已经没有力气去抗议这个篡改了。

但这个计算\textbf{确实不是"数学模型"}。它是真实的物理描述。这个体系有一个极其重大的后果：\textbf{地球不再是宇宙的中心。}人类不处在一个被上帝特别为"他的造物"安排在宇宙制高点的位置上。地球和火星、木星一样，是一颗绕着太阳运转的行星。

\subsection*{二、伽利略：望远镜对准天空}

1609年，意大利数学家\textbf{伽利略·伽利莱}听说荷兰人发明了一种"能把远处的东西放大的管子"。他没有去买这种管子——他自己做了一架更好的。然后他没有用它来看远处的船，而是把它对准了天空。他看到的东西彻底粉碎了已经主导了天文学将近一千五百年的亚里士多德和托勒密体系：

\begin{itemize}
  \item 月球表面是不光滑的——有山脉、有环形山，和地球一样。\textbf{月球不是"完美天体"}。
  \item 木星有四颗卫星绕着它转。如果木星有卫星，那这些东西就不是全部绕着地球转。\textbf{地球不是唯一的"旋转中心"}。
  \item 金星有相位变化——像月亮一样从新月变成满月。这种变化只有在金星\textbf{绕着太阳转}时才会出现。
\end{itemize}

伽利略不是一个"孤胆英雄独自对抗教会"的角色。他写《关于托勒密和哥白尼两大世界体系的对话》时，把教皇乌尔班八世的观点放在了支持地心说的角色辛普利奇奥嘴里——辛普利奇奥在意大利语里的意思是"傻瓜"。伽利略不是一个不懂政治的人，但他控制不了自己的讽刺冲动。

1633年，宗教裁判所审判了伽利略。他被迫跪在地上宣读一份悔过书，宣布放弃"地球围绕太阳转这一错误观点"。传说他在结束时低声嘟\rube{哝}{nong}了一句——"\textbf{但它确实在动}"（Eppur si muove）。这个故事很可能是杜撰的，但这句话精准地概括了他那一代科学家和教会之间的根本冲突：\textbf{大自然有它自己运行的规律——不管你嘴上承不承认}。

\subsection*{三、"我不做假设"——牛顿的上帝与机器}

1687年，英国人\textbf{艾萨克·牛顿}出版了《自然哲学的数学原理》。这本书在物理学上的突破太过惊人——牛顿用几个简洁的公式（万有引力定律、三大运动定律）统一描述了一个苹果从树上掉下和一颗行星绕着太阳运转\textbf{是同一个原因}。在此之前，人们认为"天上的规律"和"地上的规律"是两回事；在此之后，牛顿证明了——\textbf{宇宙是单一的，服从统一的物理规则}。

牛顿本人是一个极其虔诚的信徒。他生命的最后几十年投入了大量的精力去研究《圣经》中的预言和历史上的年代学，试图把上帝的创世日期精确地计算出来。但他的物理学造出了一个他从未想过要造的后果——\textbf{一个不需要上帝来持续推动的宇宙}。在他的体系里，行星不需要天使来推——引力自动把它们保持在轨道上。宇宙像一只巨大的钟表，上了发条之后就自己运行了。这个"不需要上帝日常维持的世界"给后来的思想家提供了一个深刻的启发：\textbf{如果大自然不需要不断的神意干预来运转——那人类社会呢？政治制度是不是也有它自己的、可以被发现和理解的"自然规律"？}

\subsection*{四、把科学方法搬到政治中来——启蒙运动}

17世纪末到18世纪，欧洲出现了一些开始把"分析自然的方法"应用于"分析社会""分析政治"的思想者。这场运动后来被广泛地称为\textbf{启蒙运动}。它的核心冲动不是"打倒上帝"，而是\textbf{把一切事物——政治制度、法律、道德、社会结构——都放到理性的天平上秤一秤，看看它们能不能为自己提供合理的解释}。

\begin{itemize}
  \item \textbf{约翰·洛克}（英国）：在他1689年的《政府论两篇》中论证，政府的合法性\textbf{不是来自上帝或血缘世袭}。人类生来具有生命、自由和财产的自然权利——政府是人民自愿放弃一部分权利后建立的契约产物。如果这个政府不断侵害它所服务的人民的权利——\textbf{人民有权利推翻它}。一个多世纪后，美国《独立宣言》的起草人杰斐逊几乎一字不差地照搬了这段话。
  \item \textbf{孟德斯鸠}（法国）：在1748年的《论法的精神》里分析了欧洲各国不同政体的运作机制之后，提出了一个被后人推崇备至的观点：\textbf{立法权、行政权和司法权应该分别由三个不同的人或机构执掌}，这样才能防止任何一个人把所有权力集中在一起。"\textbf{权力只有靠权力来制约}"。这套"三权分立"的框架，直接影响了美国宪法的创作者。
  \item \textbf{卢梭}（法国/日内瓦）：1762年的《社会契约论》以一句震惊了全欧洲的开场白起笔："\textbf{人生而自由，却无往不在枷锁之中。}"卢梭坚信政治权力必须来自于\textbf{全体人民的共同意志}，而不能被任何一个专制君主或少数精英垄断。他的这句话和它所承载的激进信念，后来变成了法国大革命最响亮的战斗口号。
\end{itemize}

\shihai{思想从哪里来——启蒙运动的物质基础}{
  启蒙思想不是在真空里诞生出来的。它们背后有几层坚实的物质原因：

  \begin{itemize}
    \item \textbf{印刷与识字}：古腾堡印刷术在启蒙运动之前已经传播了两百多年。到18世纪，报纸、杂志（如《旁观者》）、便宜的书籍和词典在一个识字率持续上涨的社会里形成了稳定的需求——不再只有修士和贵族才能读书写字，中产商人、手工业者也开始阅读和写作。思想有了载体。
    \item \textbf{咖啡厅和沙龙}：17世纪末到18世纪，咖啡馆在伦敦和巴黎数量密集地出现（伦敦的咖啡馆到18世纪初有数百家）。只需要花一个便士买一杯咖啡，一个人就可以坐一个下午——看报纸，听陌生人辩论政治，和一个完全不认识的人争论法律应不应该是国王说的算。咖啡馆后来被叫做"\textbf{便士大学}"——因为它用最低的门槛把知识传播给了任何负担得起一杯咖啡的人。在法国，贵族和中产阶级的沙龙成为思想碰撞的集中场所。
    \item \textbf{三十年战争和宗教战争留下的创伤}：在目睹了以信仰为名的长期互相屠杀之后，欧洲知识界普遍厌恶"为神学概念继续杀人和被杀"。一个寻找\textbf{超越宗教分歧的、基于人类共识的政治规则}的冲动，成为了启蒙运动深深的底层动力。不是这些人"更聪明"了，是他们被宗教战争的残酷和荒谬教育了。
  \end{itemize}

  所有这些力量——技术（印刷）、经济（城市消费）、社会（中产阶层崛起）、历史创伤（宗教战争后遗症）——交织在一起，\textbf{把"用脑子想事情"从一个少数神职人员的专利，变成了一个逐渐扩大的社会群体的集体行为}。美国独立和法国大革命的思想弹药，都是在这个时代开始生产和扩散的。但它们最终以什么样的方式、走什么样的道路变成了现实中枪炮、断头台和宪法——那是另一个充满了未预料后果的故事了。
}

\subsection*{五、本课要点}

\begin{center}
\begin{tabular}{|c|c|p{7cm}|}
\hline
\textbf{人物/事件} & \textbf{年代} & \textbf{核心贡献} \\
\hline
哥白尼《天体运行论》 & 1543年 & 日心说：地球围绕太阳运转 \\
\hline
伽利略用望远镜观察天空 & 1609—1633年 & 木星卫星、金星相位——证实非所有天体绕地运转 \\
\hline
牛顿《原理》 & 1687年 & 万有引力和三大运动定律——统一的宇宙 \\
\hline
洛克《政府论两篇》 & 1689年 & 政府合法性来自人民同意而非神授 \\
\hline
孟德斯鸠《论法的精神》 & 1748年 & 三权分立——用权力制约权力 \\
\hline
卢梭《社会契约论》 & 1762年 & 主权在民——"人生而自由" \\
\hline
\end{tabular}
\end{center}

\vspace{0.3cm}
{\large\textcolor{orange!80!black}{\textbf{想一想}}}
\begin{enumerate}
  \item 牛顿证明了"宇宙不需要上帝来每天推动"——但牛顿本人很虔诚。科学发现能不能和宗教信仰"和解"？还是要彻底取代？
  \item 启蒙思想家把"人为什么有权统治人"这个问题放在了理性的天平上。今天你认为什么样的政府是合法的？一个政府可以不通过选举而通过别的方式获得合法性吗？
  \item 咖啡馆在18世纪被称为"便士大学"——任何人都可以花一点钱进来听辩论、读报纸。今天的什么空间扮演了类似的角色？
\end{enumerate}

\vspace{0.3cm}
\textcolor{blue!70!black}{\textbf{延伸阅读}}
\begin{itemize}
  \item 纪录片《宇宙》——卡尔·萨根主持的经典科学史系列，前几集详细讲述了哥白尼、伽利略和牛顿的故事
  \item 《启蒙时代》——以赛亚·伯林编，了解启蒙运动的全貌
\end{itemize}

\vspace{0.5cm}
\begin{center}
\textcolor{blue!70!black}{\large —— 从查理曼到启蒙运动，欧洲七课走完一千年。政治制度不是被"设计"出来的，是在反复碰撞、妥协、失败中一层一层堆出来的。——}
\end{center}

\end{document}
